\documentclass[11pt,a4paper]{article}

\usepackage{graphicx} % Required for inserting images
\usepackage{authblk}
\usepackage[utf8]{inputenc}
\usepackage[T1]{fontenc}
\usepackage[spanish,es-nodecimaldot]{babel}
\usepackage{lmodern}
\usepackage{microtype}
\usepackage{amsmath,amssymb,amsthm,mathtools}
\usepackage{geometry}
\geometry{margin=2.5cm}
\usepackage{booktabs}
\usepackage{tabularx}
\usepackage{enumitem}
\usepackage{hyperref}
\hypersetup{colorlinks=true,linkcolor=blue,urlcolor=blue,citecolor=blue}

\theoremstyle{definition}
\newtheorem{definicion}{Definición}[section]
\newtheorem{problema}[definicion]{Problema}
\newtheorem{supuesto}[definicion]{Supuesto}
\newtheorem{observacion}[definicion]{Observación}
\newtheorem{ejemplo}[definicion]{Ejemplo}

\theoremstyle{plain}
\newtheorem{proposicion}[definicion]{Proposición}

\newcommand{\N}{\mathbb{N}}
\newcommand{\Z}{\mathbb{Z}}
\newcommand{\R}{\mathbb{R}}
\DeclareMathOperator*{\argmax}{arg\,max}

\title{Optimización del Proceso de Aprendizaje Federado}
\author{Alejandro Martínez Orozco y Jon Henry Vásquez Mesa}

\affil{Análisis y Diseño de Algoritmos}
\affil{Universidad de Caldas}

\date{6 de noviembre de 2026}

\begin{document}

\maketitle

\begin{abstract}
Este documento formaliza los problemas de \textbf{distribución de presupuesto} y \textbf{selección de participantes} en el contexto del entrenamiento de un modelo de IA de manera federada: diferentes nodos computacionales son propietarios de datos, los cuales contribuyen simultáneamente al entrenamiento del modelo. El matiz del problema a tratar surge cuando se considera la \textbf{privacidad de los datos} utilizados para el entrenamiento; con el fin de procurar su preservación, se añade ruido a los datos, el cual implica una \textbf{delegación de presupuesto} a cada nodo de tal manera que los datos en general se mantengan útiles para el aprendizaje.

Se definen matemáticamente los problemas de \textbf{distribución óptima del presupuesto de introducción de ruido} y la \textbf{selección óptima de nodos participantes} en cada etapa de entrenamiento, además de algunas variantes para ambos. Cada algoritmo se formaliza con: una \textbf{entrada de datos} y su \textbf{tamaño} correspondiente, una \textbf{salida de datos} y su \textbf{objetivo principal}. El documento no incluye todavía algoritmos ni código, puesto que su propósito es fijar el modelo, la notación y las decisiones de modelado que afectarán el análisis de complejidad posterior.
\end{abstract}

\newpage
\tableofcontents
\newpage

%=====================================================================
\section{Notación y convenciones}
%=====================================================================

\subsection{Conjuntos, parámetros y variables}

Sea $V=\{1,\dots,N\}$ el conjunto de nodos de la red. Para un subconjunto
$S\subseteq V$ se escribe $x_S=(x_i)_{i\in S}$. Si $G=(V,E)$ es un grafo, se
denota por $N(i)=\{j\in V:\{i,j\}\in E\}$ al conjunto de vecinos del nodo $i$.

Los parámetros y variables que se usan a lo largo del documento son:

\begin{description}[leftmargin=*,style=nextline]
  \item[$N$] número de nodos de la red.
  \item[$B$] presupuesto total de ruido a repartir en una ronda, en unidades
  enteras.
  \item[$B_t$] presupuesto disponible en la ronda $t$.
  \item[$P_i$] presupuesto total del nodo $i$ a lo largo de todas las rondas.
  \item[$x_i$] unidades de presupuesto asignadas al nodo $i$.
  \item[$k$] número máximo de nodos que pueden seleccionarse en una ronda.
  \item[$c_i$] costo de participación del nodo $i$.
  \item[$C$] presupuesto de participación disponible.
  \item[$T$] número de rondas.
  \item[$q_i$] probabilidad de fallo del nodo $i$ en una ronda.
  \item[$A_t$] conjunto de nodos activos en la ronda $t$.
  \item[$w_{ij}$] peso de la influencia del nodo $j$ sobre el nodo $i$.
\end{description}

\begin{definicion}[Concavidad discreta]
Una función $u:\Z_{\ge0}\to\R$ es \emph{creciente} si
$u(j+1)\ge u(j)$ para todo $j\ge0$, y es \emph{cóncava en sentido discreto}
si sus ganancias marginales
\[
  \Delta(j)=u(j)-u(j-1),\qquad j\ge1,
\]
son no crecientes. Se adopta la normalización $u(0)=0$.
\end{definicion}

\begin{definicion}[Funciones de conjunto monótonas y submodulares]
Sea $f:2^V\to\R_{\ge0}$. Se dice que $f$ es \emph{monótona} si
$A\subseteq B\subseteq V$ implica $f(A)\le f(B)$. Se dice que $f$ es
\emph{submodular} si para todo $A\subseteq B\subseteq V$ y todo
$i\in V\setminus B$ se cumple
\[
  f(A\cup\{i\})-f(A)\;\ge\; f(B\cup\{i\})-f(B).
\]
Equivalentemente, agregar un nodo a un conjunto grande aporta una ganancia
marginal menor o igual que agregarlo a uno pequeño.
\end{definicion}

\subsection{Tamaño de la entrada de datos}

\begin{definicion}[Tamaño de entrada $n$ según variaciones]\label{def:n}
Para los Problemas~\ref{prob:reparto} y~\ref{prob:seleccion} y para las
Variantes~\ref{prob:costos} y~\ref{prob:rondas}, se define $n=N$. Para la
Variante~\ref{prob:red} y para el Problema~\ref{prob:integrado}, se define
$n=N+|E|$. Los parámetros $B$, $k$, $C$, $T$ y las probabilidades $q_i$ se
reportan por separado como parámetros de la instancia. Cuando la codificación
de un parámetro numérico afecte la complejidad, se indicará explícitamente si
se usa representación unaria o binaria.
\end{definicion}

En otras palabras, lo que crece cuando la instancia se hace más grande es el
número de nodos $N$ y, en las variantes con red, también el número de aristas
$|E|$. Los parámetros numéricos se mantienen fijos o se reportan aparte.

\subsection{Modelo de acceso a la entrada}

\begin{supuesto}[Acceso por oráculo]\label{sup:oraculo}
Las funciones de utilidad $u_i$ se representan mediante un oráculo que
devuelve $u_i(j)$ o su ganancia marginal $\Delta_i(j)$ en tiempo $O(1)$. La
función de conjunto $f$ se representa mediante un oráculo que devuelve $f(S)$
en tiempo $O(1)$ para cualquier $S\subseteq V$. El costo de las consultas al
oráculo se contabiliza en la función de eficiencia del algoritmo.
\end{supuesto}

\begin{observacion}[Representación explícita]
\label{obs:representacion}
Si las utilidades se entregaran como una tabla con las $B$ ganancias
marginales de cada nodo, la entrada tendría $\Theta(NB)$ números. Si $f$ se
entregara de forma explícita, su descripción podría tener tamaño exponencial en
$N$. El modelo de oráculo del Supuesto~\ref{sup:oraculo} es el que se adopta en
este documento y es el que permite que $n=N$ (o $n=N+|E|$) sea una medida
combinatoria significativa del tamaño de la instancia.
\end{observacion}

\subsection{Supuestos generales}

\begin{supuesto}[Supuestos del modelo]\label{sup:modelo}
Salvo que se indique lo contrario:
\begin{enumerate}[leftmargin=*,label=\arabic*)]
  \item Los presupuestos se miden en unidades enteras no negativas.
  \item Las utilidades son no negativas, crecientes y cóncavas en sentido
  discreto, con $u_i(0)=0$.
  \item La función de conjunto $f$ es no negativa, monótona y submodular.
  \item La red no tiene autoaristas y los pesos de influencia son no negativos.
  \item Los nodos y las aristas forman parte de la entrada; las utilidades y
  $f$ se acceden por oráculo.
\end{enumerate}
\end{supuesto}

%=====================================================================
\section{Problema 1: reparto óptimo del presupuesto}
%=====================================================================

\begin{problema}\label{prob:reparto} Reparto óptimo del presupuesto

\textbf{Entrada.} Un entero $N\in\N$, un presupuesto $B\in\Z_{\ge0}$ y, para
cada nodo $i\in V$, una función de utilidad $u_i:\Z_{\ge0}\to\R_{\ge0}$
creciente y cóncava en sentido discreto, con $u_i(0)=0$, accesible por oráculo
según el Supuesto~\ref{sup:oraculo}.

\textbf{Salida.} Un vector de asignación
$x=(x_1,\dots,x_N)\in\Z_{\ge0}^N$.

\textbf{Factibilidad.} $\sum_{i=1}^{N} x_i\le B$.

\textbf{Objetivo.} Maximizar la utilidad total
\[
  \sum_{i=1}^{N} u_i(x_i).
\]

\textbf{Tamaño de la entrada.} $n=N$, de acuerdo con la
Definición~\ref{def:n}. El presupuesto $B$ se reporta como parámetro adicional.
\end{problema}

\begin{observacion}[Uso completo del presupuesto]
Como las funciones $u_i$ son no decrecientes, toda solución factible con
$\sum_i x_i<B$ puede completarse asignando las unidades restantes a cualquier
nodo sin disminuir la utilidad total. Por lo tanto, sin pérdida de
generalidad, puede exigirse $\sum_i x_i=B$.
\end{observacion}

\begin{observacion}[Codificación de $B$ y complejidad]
El algoritmo voraz que se estudiará más adelante asigna una unidad de
presupuesto a la vez y tiene costo $O(B\log N)$ usando una cola de prioridad.
Este costo es polinomial en $B$ y en $N$. Si $B$ se codificara en binario, el
algoritmo sería pseudo-polinomial en el tamaño de la representación de $B$. El
enunciado del curso adopta el parámetro $B$ como parte de la instancia y
reporta la complejidad como $T(N,B)=O(B\log N)$; esta distinción se deja
explícita para que la definición de $n$ no sea ambigua.
\end{observacion}

%=====================================================================
\section{Problema 2: selección de nodos}
%=====================================================================

\begin{problema}\label{prob:seleccion} Selección con cupo fijo

\textbf{Entrada.} Un entero $N\in\N$, un entero $k$ con $0\le k\le N$ y una
función de conjunto $f:2^V\to\R_{\ge0}$ monótona y submodular, accesible por
oráculo según el Supuesto~\ref{sup:oraculo}.

\textbf{Salida.} Un subconjunto $S\subseteq V$ con $|S|=k$.

\textbf{Objetivo.} Maximizar $f(S)$.

\textbf{Tamaño de la entrada.} $n=N$, de acuerdo con la
Definición~\ref{def:n}. El cupo $k$ se reporta como parámetro adicional.
\end{problema}

\begin{observacion}[Casos límite]
Si $k=0$, la única solución factible es $S=\emptyset$ y el problema es trivial.
Si $k>N$, no existe solución factible; en ese caso la instancia se declara
infactible. Si $k=N$, la única solución es $S=V$.
\end{observacion}

\begin{observacion}[Relación con la Parte 1]
El Problema~\ref{prob:seleccion} no incluye el reparto del presupuesto: una vez
elegido $S$, la asignación $\{x_i\}_{i\in S}$ se resolvería con el
Problema~\ref{prob:reparto}. Los dos problemas se combinan en el
Problema~\ref{prob:integrado}.
\end{observacion}

%=====================================================================
\section{Variante 1: costos de participación}
%=====================================================================

\begin{problema}\label{prob:costos} Selección con costos y presupuesto

\textbf{Entrada.} Un entero $N\in\N$, un costo $c_i>0$ para cada $i\in V$, un
presupuesto de participación $C\in\Z_{\ge0}$ y una función de conjunto
$f:2^V\to\R_{\ge0}$ monótona y submodular, accesible por oráculo.

\textbf{Salida.} Un subconjunto $S\subseteq V$.

\textbf{Factibilidad.} $\sum_{i\in S} c_i\le C$.

\textbf{Objetivo.} Maximizar $f(S)$.

\textbf{Tamaño de la entrada.} $n=N$, de acuerdo con la
Definición~\ref{def:n}. Los costos $c_i$ forman parte de la entrada y $C$ se
reporta como parámetro adicional.
\end{problema}

\begin{observacion}[El caso base es un caso particular]
Si $c_i=1$ para todo $i\in V$ y $C=k$, la restricción
$\sum_{i\in S} c_i\le C$ equivale a $|S|\le k$. Como $f$ es monótona, toda
solución óptima puede extenderse hasta tener exactamente $k$ elementos (si
$k\le N$), recuperándose el Problema~\ref{prob:seleccion}.
\end{observacion}

\begin{observacion}[Cota de aproximación]
La cota clásica $(1-1/e)$ para el algoritmo voraz del Problema~\ref{prob:seleccion}
no se aplica directamente a esta variante. La versión correcta para cobertura
con presupuesto (\emph{budgeted maximum coverage}) debe citarse o demostrarse
en la etapa correspondiente del proyecto.
\end{observacion}

%=====================================================================
\section{Variante 2: $T$ rondas con fallos}
%=====================================================================

\subsection{Proceso de fallos}

\begin{definicion}[Proceso de fallos]
Sea $(\Omega,\mathcal{F},\Pr)$ un espacio de probabilidad y sea
$(\mathcal{F}_t)_{t=0}^{T}$ una filtración. Un \emph{proceso de fallos} es una
sucesión de conjuntos activos $(A_t)_{t=1}^{T}$ tal que $A_t\in\mathcal{F}_t$ y
$A_t\subseteq V$ para cada $t$. Se consideran dos instancias:
\begin{enumerate}[leftmargin=*,label=\roman*)]
  \item \emph{Fallos independientes por ronda:}
  $A_t=\{i\in V: F_{i,t}=0\}$, donde
  $F_{i,t}\sim\mathrm{Bernoulli}(q_i)$ son independientes para todo $i$ y $t$.
  \item \emph{Fallos absorbentes:} $A_0=V$ y
  $A_t=A_{t-1}\setminus\{i\in A_{t-1}: F_{i,t}=1\}$, de modo que un nodo que
  falla no vuelve a estar activo.
\end{enumerate}
\end{definicion}

\subsection{Problema con rondas}

\begin{problema}\label{prob:rondas} Selección y asignación en $T$ rondas con fallos

\textbf{Entrada.} Un entero $N\in\N$, un número de rondas $T\in\N$, un proceso
de fallos $(A_t)_{t=1}^{T}$ con probabilidades $q_i$, presupuestos por ronda
$B_t\in\Z_{\ge0}$, presupuestos totales por nodo $P_i\in\Z_{\ge0}$, y para cada
ronda $t$ una función de conjunto $f_t:2^V\to\R_{\ge0}$ monótona y submodular y
utilidades $u_i^t:\Z_{\ge0}\to\R_{\ge0}$ crecientes y cóncavas.

\textbf{Decisión.} Una política $\pi=(\pi_1,\dots,\pi_T)$. En la ronda $t$, la
política observa el historial
\[
  H_{t-1}=\sigma\bigl(A_1,S_1,x_1,\dots,A_{t-1},S_{t-1},x_{t-1}\bigr)
\]
y el conjunto activo $A_t$, y elige
\[
  (S_t,x_t)=\pi_t(H_{t-1},A_t),
\]
donde $S_t\subseteq A_t$ y $x_t=(x_i^t)_{i\in S_t}\in\Z_{\ge0}^{S_t}$.

\textbf{Factibilidad.} Para cada ronda $t$,
\[
  \sum_{i\in S_t} x_i^t\le B_t,
  \qquad
  \sum_{t=1}^{T} x_i^t\le P_i\quad\text{para todo }i\in V.
\]

\textbf{Salida.} Una política factible $\pi$.

\textbf{Objetivo.} Maximizar la utilidad total esperada
\[
  \max_{\pi}\;
  \mathbb{E}\left[
    \sum_{t=1}^{T}
    \left(
      f_t(S_t)+\sum_{i\in S_t} u_i^t(x_i^t)
    \right)
  \right].
\]

\textbf{Tamaño de la entrada.} $n=N$, de acuerdo con la
Definición~\ref{def:n}. El número de rondas $T$ se reporta como parámetro
adicional.
\end{problema}

\begin{observacion}[Momento de la decisión]
Se supone que el conjunto activo $A_t$ se observa antes de elegir $S_t$. Si los
fallos ocurrieran después de la decisión, la política debería ser no
anticipativa respecto de $A_t$ y el modelo cambiaría. Esta convención debe
confirmarse con el enunciado del curso.
\end{observacion}

\begin{observacion}[Esperanza versus peor caso]
El objetivo del Problema~\ref{prob:rondas} es la utilidad esperada. Una
alternativa es optimizar el peor caso sobre las trayectorias de fallos. Ambas
versiones son válidas como formalización, pero producen análisis de
aproximación distintos; la elección debe quedar explícita.
\end{observacion}

\begin{observacion}[Caso estático]
Si $T=1$ y $A_1=V$, el Problema~\ref{prob:rondas} se reduce a la combinación
de una selección con cupo o costo y un reparto de presupuesto en una sola
ronda.
\end{observacion}

%=====================================================================
\section{Variante 3: interacciones en la red}
%=====================================================================

\subsection{Modelo general}

\begin{problema}\label{prob:red} Reparto con externalidades de red

\textbf{Entrada.} Un entero $N\in\N$, un grafo no dirigido $G=(V,E)$ y pesos de
influencia $w_{ij}\ge0$ para cada par ordenado $(i,j)$ con $\{i,j\}\in E$, un
presupuesto $B\in\Z_{\ge0}$ y, para cada nodo $i\in V$, una función de utilidad
\[
  u_i:\Z_{\ge0}^{1+|N(i)|}\to\R_{\ge0},
  \qquad
  u_i(x_i,x_{N(i)}),
\]
creciente y cóncava en cada una de sus coordenadas, con $u_i(0,\dots,0)=0$.

\textbf{Salida.} Un vector de asignación
$x=(x_1,\dots,x_N)\in\Z_{\ge0}^N$.

\textbf{Factibilidad.} $\sum_{i=1}^{N} x_i\le B$.

\textbf{Objetivo.} Maximizar la utilidad total de la red
\[
  \sum_{i=1}^{N} u_i(x_i,x_{N(i)}).
\]

\textbf{Tamaño de la entrada.} $n=N+|E|$, de acuerdo con la
Definición~\ref{def:n}. El presupuesto $B$ se reporta como parámetro adicional.
\end{problema}

\subsection{Instancia separable}

\begin{ejemplo}[Externalidades separables]\label{ej:red}
Una forma concreta y tratable de modelar la influencia es
\[
  u_i(x_i,x_{N(i)})
  =
  v_i(x_i)
  +
  \sum_{j\in N(i)} w_{ij}\,g_{ij}(x_j),
\]
donde $v_i:\Z_{\ge0}\to\R_{\ge0}$ y $g_{ij}:\Z_{\ge0}\to\R_{\ge0}$ son
crecientes y cóncavas en sentido discreto, $v_i(0)=0$ y $g_{ij}(0)=0$. En este
caso, la utilidad total es
\[
  \sum_{i\in V} v_i(x_i)
  +
  \sum_{\{i,j\}\in E} \bigl(w_{ij}\,g_{ij}(x_j)+w_{ji}\,g_{ji}(x_i)\bigr),
\]
que sigue siendo una suma de funciones cóncavas y por lo tanto conserva la
concavidad del problema de asignación.
\end{ejemplo}

\begin{observacion}[Efecto sobre la selección]
Si la utilidad de seleccionar un nodo depende de los vecinos seleccionados,
la función de conjunto resultante $f_G(S)$ no es necesariamente submodular.
La submodularidad debe imponerse como hipótesis o demostrarse para la
instancia concreta; no se hereda automáticamente de las utilidades locales.
\end{observacion}

%=====================================================================
\section{Problema integrado}
%=====================================================================

\begin{problema}\label{prob:integrado} Problema integrado

\textbf{Entrada.} Un entero $N\in\N$, un número de rondas $T\in\N$, un grafo
$G=(V,E)$ con pesos $w_{ij}\ge0$, un proceso de fallos $(A_t)_{t=1}^{T}$ con
probabilidades $q_i$, presupuestos por ronda $B_t$, presupuestos totales por
nodo $P_i$, cupos $k_t$ (o costos $c_i$ y presupuestos $C_t$), y para cada
ronda $t$ funciones $f_t:2^V\to\R_{\ge0}$ monótonas y submodulares y utilidades
$u_i^t$ crecientes y cóncavas.

\textbf{Decisión.} Una política $\pi=(\pi_1,\dots,\pi_T)$ que en cada ronda
elige $S_t\subseteq A_t$ y una asignación $x_t=(x_i^t)_{i\in S_t}$.

\textbf{Factibilidad.} Para cada ronda $t$:
\[
  S_t\subseteq A_t,
  \qquad
  |S_t|\le k_t
  \quad\text{(o bien }\textstyle\sum_{i\in S_t}c_i\le C_t\text{)},
\]
\[
  \sum_{i\in S_t} x_i^t\le B_t,
  \qquad
  \sum_{t=1}^{T} x_i^t\le P_i\quad\text{para todo }i\in V.
\]

\textbf{Salida.} Una política factible $\pi$.

\textbf{Objetivo.} Maximizar
\[
  \max_{\pi}\;
  \mathbb{E}\left[
    \sum_{t=1}^{T}
    \left(
      f_t(S_t)
      +
      \sum_{i\in S_t} u_i^t\bigl(x_i^t,x_{N(i)}^t\bigr)
    \right)
  \right].
\]

\textbf{Tamaño de la entrada.} $n=N+|E|$, de acuerdo con la
Definición~\ref{def:n}. Los parámetros $T$, $B_t$, $P_i$, $k_t$ y $C_t$ se
reportan por separado.
\end{problema}

\begin{observacion}[Alcance de las externalidades]
En el objetivo del Problema~\ref{prob:integrado} se contabiliza la utilidad de
los nodos seleccionados. Si se desea que el presupuesto asignado a un nodo
afecte también la utilidad de vecinos no seleccionados, se puede reemplazar el
índice $i\in S_t$ por $i\in A_t$ y definir $x_i^t=0$ para $i\notin S_t$. Esta
elección debe quedar explícita en el informe.
\end{observacion}

\begin{observacion}[Los problemas anteriores son subproblemas]
El Problema~\ref{prob:reparto} se obtiene del Problema~\ref{prob:integrado}
fijando $T=1$, sin fallos, sin red, y fijando previamente el conjunto de nodos
participantes. El Problema~\ref{prob:seleccion} se obtiene fijando $T=1$, sin
fallos, sin red, y sin reparto de presupuesto. Las variantes agregan,
respectivamente, costos de participación, múltiples rondas con fallos e
interacciones de red.
\end{observacion}

\subsection{Resumen}

\begin{table}[htbp]
\centering
\small
\renewcommand{\arraystretch}{1.25}
\begin{tabularx}{\textwidth}{@{}p{2.6cm}p{1.7cm}X X X@{}}
\toprule
Problema & $n$ & Entrada & Salida & Objetivo \\
\midrule
\ref{prob:reparto} Reparto &
$N$ &
$N,B,u_i$ &
$x\in\Z_{\ge0}^N$ &
$\max\sum_i u_i(x_i)$ \\
\addlinespace
\ref{prob:seleccion} Selección &
$N$ &
$N,k,f$ &
$S\subseteq V,\ |S|=k$ &
$\max f(S)$ \\
\addlinespace
\ref{prob:costos} Costos &
$N$ &
$N,c_i,C,f$ &
$S\subseteq V$ &
$\max f(S)$ s.t. $\sum_{i\in S}c_i\le C$ \\
\addlinespace
\ref{prob:rondas} Rondas &
$N$ &
$N,T,q_i,B_t,P_i,f_t,u_i^t$ &
política $\pi$ &
$\max_\pi\mathbb{E}\bigl[\sum_t f_t(S_t)+\sum_{i\in S_t}u_i^t(x_i^t)\bigr]$ \\
\addlinespace
\ref{prob:red} Red &
$N+|E|$ &
$N,G,w_{ij},B,u_i$ &
$x\in\Z_{\ge0}^N$ &
$\max\sum_i u_i(x_i,x_{N(i)})$ \\
\addlinespace
\ref{prob:integrado} Integrado &
$N+|E|$ &
todos los anteriores &
política $\pi$ &
$\max_\pi\mathbb{E}\bigl[\sum_t f_t(S_t)+\sum_{i\in S_t}u_i^t(x_i^t,x_{N(i)}^t)\bigr]$ \\
\bottomrule
\end{tabularx}
\caption{Resumen de los problemas formalizados.}
\label{tab:resumen}
\end{table}

%=====================================================================
\section{Supuestos, casos límite y decisiones abiertas}
%=====================================================================

\subsection{Casos límite}

\begin{itemize}[leftmargin=*,label=$\bullet$]
  \item $k=0$: la selección es vacía y el Problema~\ref{prob:seleccion} es
  trivial.
  \item $k>N$: la instancia del Problema~\ref{prob:seleccion} es infactible.
  \item $B=0$ o $C=0$: no se asigna presupuesto o no se selecciona ningún nodo,
  respectivamente.
  \item $P_i=0$: el nodo $i$ no puede participar en ninguna ronda.
  \item $N=1$: los problemas de selección son triviales y el reparto asigna
  todo el presupuesto al único nodo.
  \item Nodos aislados en la Variante~\ref{prob:red}: $N(i)=\emptyset$ y la
  utilidad depende solo de $x_i$.
  \item $q_i=0$: el nodo nunca falla; $q_i=1$: el nodo nunca está activo.
  \item $T=1$: el Problema~\ref{prob:rondas} se reduce a una sola ronda.
\end{itemize}

\subsection{Decisiones de modelado por confirmar}

\begin{enumerate}[leftmargin=*,label=\arabic*)]
  \item \textbf{Unidades de presupuesto.} Se asume que el presupuesto es
  discreto y se asigna en unidades enteras. Si el presupuesto fuera continuo,
  el algoritmo voraz cambiaría por un esquema de reparto por umbrales y la
  función de eficiencia sería distinta.
  \item \textbf{Representación de las utilidades.} Se adopta el modelo de
  oráculo del Supuesto~\ref{sup:oraculo}. Si el curso exige una representación
  explícita, el tamaño de la entrada y la complejidad deben ajustarse según la
  Observación~\ref{obs:representacion}.
  \item \textbf{Momento de los fallos.} Se asume que $A_t$ se observa antes de
  decidir en la ronda $t$. Si los fallos ocurren después, la política debe ser
  no anticipativa y el análisis cambia.
  \item \textbf{Objetivo en varias rondas.} Se adopta la utilidad esperada.
  La alternativa de peor caso debe indicarse explícitamente si se prefiere.
  \item \textbf{Presupuesto entre rondas.} Se asume que cada nodo tiene un
  presupuesto total $P_i$ y que el presupuesto no gastado no se transfiere
  entre nodos. Si se permite transferencia, debe agregarse esa restricción al
  modelo.
  \item \textbf{Interacciones de red.} La forma separable del
  Ejemplo~\ref{ej:red} preserva la concavidad. Si se usa una utilidad general
  $u_i(x_i,x_{N(i)})$, deben verificarse las hipótesis de concavidad y
  submodularidad que requiera cada algoritmo.
\end{enumerate}

%=====================================================================
\section{Lista de verificación de la rúbrica}
%=====================================================================

\begin{itemize}[leftmargin=*,label=$\square$]
  \item Se define con precisión el tamaño de la entrada $n$ para cada problema
  y se indica qué crece cuando la instancia se hace más grande.
  \item Cada problema especifica entrada, salida, restricciones de factibilidad
  y objetivo.
  \item Se distingue entre los problemas base y las tres variantes obligatorias.
  \item Se explicita el modelo de acceso a las utilidades y a la función de
  conjunto.
  \item Se dejan identificadas las decisiones de modelado que afectan el
  análisis de complejidad y las cotas de aproximación.
\end{itemize}

\end{document}